\documentclass[aspectratio=169,10pt]{beamer}
\usetheme[numbering=fraction,progressbar=frametitle]{metropolis}
\usepackage{fontspec}
\setsansfont{Latin Modern Sans}
\usepackage{amsmath,amssymb,mathtools,graphicx,booktabs}
\usepackage{polyglossia}
\setdefaultlanguage{french}
\setbeamertemplate{navigation symbols}{}
\setbeamerfont{frametitle}{size=\large}
\setbeamerfont{normal text}{size=\normalsize}
\setlength{\parskip}{.45em}
\hypersetup{pdfauthor={},pdfsubject={Lecture slides},pdfkeywords={},pdftitle={Probabilités : conditionnement et arbres pondérés}}
\providecommand{\N}{\mathbb{N}}
\providecommand{\Z}{\mathbb{Z}}
\providecommand{\D}{\mathbb{D}}
\providecommand{\Q}{\mathbb{Q}}
\providecommand{\R}{\mathbb{R}}
\providecommand{\C}{\mathbb{C}}
\providecommand{\vect}[1]{\overrightarrow{#1}}
\providecommand{\vi}{\vec{\imath}}
\providecommand{\vj}{\vec{\jmath}}
\providecommand{\vk}{\vec{k}}
\providecommand{\norm}[1]{\left\lVert #1\right\rVert}
\providecommand{\vc}[2]{\begin{pmatrix}#1\\#2\end{pmatrix}}
\providecommand{\deter}[2]{\det\left(#1,#2\right)}
\providecommand{\abs}[1]{\left\lvert #1\right\rvert}
\providecommand{\dd}{\,\mathrm{d}}
\providecommand{\e}{\mathrm{e}}
\providecommand{\Card}{\operatorname{Card}}
\providecommand{\E}{\operatorname{E}}
\providecommand{\V}{\operatorname{V}}
\providecommand{\degre}{^{\circ}}
\providecommand{\un}[1]{\,\mathrm{#1}}
\providecommand{\sci}[2]{#1\times10^{#2}}
\providecommand{\pgcd}{\operatorname{PGCD}}
\providecommand{\ppcm}{\operatorname{PPCM}}
\title{Probabilités : conditionnement et arbres pondérés}
\subtitle{Cours et exemples}
\date{}
\author{}
\begin{document}
\begin{frame}[plain]
\vfill
{\usebeamerfont{title}\inserttitle\par}
\medskip
{\usebeamerfont{subtitle}\insertsubtitle\par}
\medskip
{\usebeamercolor[fg]{progress bar}\rule{\textwidth}{.5pt}}
\vfill
\end{frame}
\begin{frame}{Univers, issues et événements}
Une expérience aléatoire possède plusieurs résultats possibles sans que l'on sache à l'avance lequel sera obtenu. Une issue est un résultat possible ; l'univers \(\Omega\) est leur ensemble. Un événement est une partie de \(\Omega\) ; un événement élémentaire contient une seule issue.
L'événement impossible est \(\varnothing\) ; l'événement certain est \(\Omega\).
\end{frame}
\begin{frame}{Univers, issues et événements}
\(\overline A\) est l'événement contraire de \(A\). \(A\cap B\) signifie « \(A\) et \(B\) » ; \(A\cup B\), « \(A\) ou \(B\) », éventuellement les deux. Deux événements sont incompatibles si \(A\cap B=\varnothing\) ; sinon, ils sont compatibles.
\end{frame}
\begin{frame}{Univers, issues et événements}
On lance un dé à six faces. Écrire \(\Omega\), \(A\) : « obtenir un nombre pair », \(B\) : « obtenir \(4\) », puis \(\overline A\), \(A\cap B\) et \(A\cup B\). \(A,B\) sont-ils compatibles ?
\end{frame}
\begin{frame}{Correction : Univers, issues et événements}
\(\Omega=\{1;2;3;4;5;6\}\), \(A=\{2;4;6\}\), \(B=\{4\}\), \(\overline A=\{1;3;5\}\), \(A\cap B=\{4\}\), \(A\cup B=A\). Les événements sont compatibles car ils partagent l'issue \(4\).
\end{frame}
\begin{frame}{Univers, issues et événements}
Avec \(A=\{a;b;c;d;e\}\) et \(B=\{b;e;f;g;h\}\), donner l'intersection et la réunion. Peut-on déterminer \(\overline A\) sans préciser l'univers ?
\end{frame}
\begin{frame}{Correction : Univers, issues et événements}
\(A\cap B=\{b;e\}\) et \(A\cup B=\{a;b;c;d;e;f;g;h\}\). Le complémentaire dépend de l'univers ; on ne peut pas le déterminer sans cette information.
\end{frame}
\begin{frame}{Modèle probabiliste et équiprobabilité}
À chaque issue, un modèle associe une probabilité positive ou nulle ; la somme vaut \(1\). La probabilité d'un événement est la somme des probabilités de ses issues. On a
\[P(\Omega)=1,\quad P(\varnothing)=0,\quad P(\overline A)=1-P(A).\]
En équiprobabilité, toutes les issues ont même probabilité et
\[P(A)=\dfrac{\Card(A)}{\Card(\Omega)}.\]
L'équiprobabilité est une hypothèse du modèle, pas une conséquence de la seule liste des issues.
\end{frame}
\begin{frame}{Modèle probabiliste et équiprobabilité}
On suppose le dé équilibré. Pour \(A=\) « pair » et \(C=\) « supérieur ou égal à \(5\) », calculer \(P(A)\), \(P(C)\), \(P(A\cap C)\) et \(P(A\cup C)\).
\end{frame}
\begin{frame}{Correction : Modèle probabiliste et équiprobabilité}
\(P(A)=3/6=1/2\), \(P(C)=2/6=1/3\), \(A\cap C=\{6\}\) donc \(P(A\cap C)=1/6\). \(A\cup C=\{2;4;5;6\}\), donc \(P(A\cup C)=4/6=2/3\). L'issue \(6\) ne doit pas être comptée deux fois.
\end{frame}
\begin{frame}{Modèle probabiliste et équiprobabilité}
Une punaise peut tomber pointe vers le haut ou vers le bas. Peut-on en déduire que chaque issue a probabilité \(1/2\) ? Comment proposer un modèle ?
\end{frame}
\begin{frame}{Correction : Modèle probabiliste et équiprobabilité}
Non : deux issues ne suffisent pas à imposer l'équiprobabilité. On peut répéter les lancers dans des conditions comparables et utiliser les fréquences observées pour proposer un modèle, en gardant à l'esprit les fluctuations et les caractéristiques de l'objet.
\end{frame}
\begin{frame}{Probabilité conditionnelle à partir d'un tableau}
Pour \(P(A)>0\), la probabilité de \(B\) sachant \(A\) est
\[P_A(B)=\dfrac{P(A\cap B)}{P(A)},\qquad P(A\cap B)=P(A)P_A(B).\]
Dans un tirage équiprobable au sein d'une population finie :
\[P_A(B)=\dfrac{\Card(A\cap B)}{\Card(A)}.\]
On note aussi \(P(B\mid A)\). \(P_A(B)\) et \(P_B(A)\) ne sont généralement pas égales.
\end{frame}
\begin{frame}{Probabilité conditionnelle à partir d'un tableau}
On tire uniformément un élève de la classe de \(32\) élèves étudiée précédemment. \(G\) : « garçon », \(N\) : « suit NSI ».
\end{frame}
\begin{frame}{Correction : Probabilité conditionnelle à partir d'un tableau}
\(P(G)=18/32=9/16\), \(P(N)=21/32\), \(P(G\cap N)=12/32=3/8\).
\(P_N(G)=12/21=4/7\), tandis que \(P_G(N)=12/18=2/3\).
Par la formule : \(P_N(G)=(12/32)/(21/32)=12/21\). Sachant que l'élève suit NSI, la population de référence est réduite aux \(21\) élèves concernés.
\end{frame}
\begin{frame}{Probabilité conditionnelle à partir d'un tableau}
Le conditionnement exige une probabilité non nulle de l'événement connu ; on ne divise pas par zéro.
\end{frame}
\begin{frame}{Arbres pondérés et probabilité d'un chemin}
\begin{columns}[T,onlytextwidth]
\begin{column}{.57\textwidth}
À chaque nœud, la somme des probabilités des branches sortantes vaut \(1\). Une branche du deuxième niveau porte une probabilité conditionnelle. Un chemin représente une intersection d'événements ; sa probabilité est le produit des probabilités inscrites sur ce chemin.
\end{column}
\begin{column}{.4\textwidth}
\centering
\includegraphics[width=\linewidth,height=4.5cm,keepaspectratio]{figures/seconde-14-1.pdf}
\end{column}
\end{columns}
\end{frame}
\begin{frame}{Arbres pondérés et probabilité d'un chemin}
On sait \(P(A)=0{,}6\), \(P_A(B)=0{,}7\) et \(P_{\overline A}(\overline B)=0{,}9\). Compléter les autres branches et calculer les probabilités des quatre chemins.
\end{frame}
\begin{frame}{Correction : Arbres pondérés et probabilité d'un chemin}
\(P(\overline A)=0{,}4\), \(P_A(\overline B)=0{,}3\), \(P_{\overline A}(B)=0{,}1\).
\[P(A\cap B)=0{,}6\times0{,}7=0{,}42,\quad P(A\cap\overline B)=0{,}18,\]
\[P(\overline A\cap B)=0{,}04,\qquad P(\overline A\cap\overline B)=0{,}36.\]
La somme des quatre probabilités vaut \(1\), ce qui contrôle la cohérence du modèle.
\end{frame}
\begin{frame}{Arbres pondérés et probabilité d'un chemin}
Justifier la multiplication sur le chemin \(A\) puis \(B\) à partir de la définition de \(P_A(B)\).
\end{frame}
\begin{frame}{Correction : Arbres pondérés et probabilité d'un chemin}
Comme \(P_A(B)=P(A\cap B)/P(A)\) et \(P(A)>0\), multiplier par \(P(A)\) donne \(P(A)P_A(B)=P(A\cap B)\).
\end{frame}
\begin{frame}{Passer d'un tableau à un arbre}
Pour la classe de \(32\) élèves, construire un arbre commençant par NSI / pas de NSI, puis garçon / fille. Construire ensuite l'arbre dans l'ordre inverse. Vérifier, dans les deux arbres, la probabilité de choisir une fille suivant NSI.
\end{frame}
\begin{frame}{Correction : Passer d'un tableau à un arbre}
Premier arbre : branches initiales \(P(N)=21/32\), \(P(\overline N)=11/32\). Après \(N\) : \(P_N(G)=12/21=4/7\), \(P_N(\overline G)=9/21=3/7\). Après \(\overline N\) : \(6/11\) et \(5/11\).
Second arbre : branches initiales \(P(G)=18/32=9/16\), \(P(\overline G)=14/32=7/16\). Après \(G\) : \(P_G(N)=12/18=2/3\), puis \(1/3\). Après \(\overline G\) : \(P_{\overline G}(N)=9/14\), puis \(5/14\).
Le chemin recherché vaut \((21/32)(9/21)=9/32\) dans le premier arbre et \((14/32)(9/14)=9/32\) dans le second. Le changement d'ordre conserve les probabilités des intersections mais change les probabilités des branches.
\end{frame}
\begin{frame}{Passer d'un tableau à un arbre}
Une urne contient \(3\) boules rouges et \(2\) bleues. On tire successivement deux boules \textbf{sans remise}. Construire l'arbre des couleurs et calculer la probabilité de tirer deux rouges.
\end{frame}
\begin{frame}{Correction : Passer d'un tableau à un arbre}
Premier niveau : \(R\) avec \(3/5\), \(B\) avec \(2/5\). Après \(R\), il reste \(2\) rouges et \(2\) bleues : branches \(1/2,1/2\). Après \(B\), il reste \(3\) rouges et \(1\) bleue : branches \(3/4,1/4\). Donc \(P(R_1\cap R_2)=(3/5)(1/2)=3/10\). Utiliser \(3/5\) sur la seconde branche rouge ignorerait l'absence de remise.
\end{frame}
\begin{frame}{Inversion des conditionnements et faux positifs}
Un lot contient \(10000\) composants, dont \(100\) défectueux. Un test est positif pour \(99\) des composants défectueux et pour \(495\) des composants non défectueux. On choisit uniformément un composant. \(D\) : « défectueux » ; \(T\) : « test positif ».

\begin{enumerate}
\setcounter{enumi}{0}\item Compléter le tableau croisé test / état réel.
\item Calculer \(P_D(T)\) et \(P_T(D)\).
\item Donner le nombre de faux positifs et de faux négatifs.
\item Calculer la sensibilité \(P_D(T)\) et la spécificité \(P_{\overline D}(\overline T)\).

\end{enumerate}
\end{frame}
\begin{frame}{Correction : Inversion des conditionnements et faux positifs}
center tabularc|rr|r\&\(D\)\&\(\overline D\)\&Total

 \(T\)\&99\&495\&594

\(\overline T\)\&1\&9405\&9406

 Total\&100\&9900\&10000tabularcenter
\(P_D(T)=99/100=99\,\%\). En revanche \(P_T(D)=99/594=1/6\approx16{,}67\,\%\). Un test positif n'implique donc pas ici une probabilité de défaut de \(99\,\%\).
Les \(495\) faux positifs sont les composants sains testés positifs ; le faux négatif est le composant défectueux testé négatif. La sensibilité vaut \(99\,\%\) ; la spécificité \(9405/9900=95\,\%\).
La faible proportion initiale de composants défectueux explique pourquoi les faux positifs sont plus nombreux que les vrais positifs malgré la bonne sensibilité.
\end{frame}
\begin{frame}{Inversion des conditionnements et faux positifs}
« Positif parmi les défectueux » et « défectueux parmi les positifs » correspondent à deux conditionnements différents. Toujours nommer le groupe de référence avant le calcul.
\end{frame}
\begin{frame}{Simulation et loi des grands nombres}
Lorsqu'on répète un grand nombre de fois une expérience dans le même modèle avec des répétitions indépendantes, la fréquence observée d'un événement est, sauf exception, proche de sa probabilité. Elle fluctue : elle ne s'en rapproche pas nécessairement à chaque nouvelle expérience.
\end{frame}
\begin{frame}[fragile]{Python}
\small
\begin{verbatim}
from random import Random

def frequences_six(n, graine=1):
    # n entier >= 1 ; modele de lancers independants
    rng = Random(graine)
    succes = 0
    frequences = []
    for k in range(1, n+1):
        if rng.randint(1, 6) == 6:
            succes = succes+1
        frequences.append(succes/k)
    return frequences
\end{verbatim}
\end{frame}
\begin{frame}{Simulation et loi des grands nombres}
Tracer les fréquences pour \(n=1000\), avec plusieurs graines. Comparer à \(1/6\). Pourquoi fixer la graine ? Peut-on exiger que la fréquence finale soit exactement \(1/6\) ?
\end{frame}
\begin{frame}{Correction : Simulation et loi des grands nombres}
Les tracés dépendent de la graine et fluctuent autour de \(1/6\). Fixer la graine rend la simulation reproductible. Avec \(1000\) lancers, une fréquence finale exactement \(1/6\) est même impossible, car \(1000/6\) n'est pas entier. Une fréquence proche est cohérente avec le modèle, mais n'en constitue pas une démonstration.
\end{frame}
\begin{frame}{Simulation et loi des grands nombres}
Modifier le programme pour suivre la fréquence de l'événement « résultat pair ». Quelle probabilité de référence utiliser ?
\end{frame}
\begin{frame}{Correction : Simulation et loi des grands nombres}
Remplacer le test par \texttt{rng.randint(1, 6) \% 2 == 0}. La probabilité de référence est \(3/6=1/2\). Chaque expérience doit effectuer un seul tirage ; ne pas appeler plusieurs fois le générateur pour tester un même résultat.
\end{frame}
\end{document}
